\section{Метод наименьших квадратов}
\label{sec:least-squares}

\subsection{Линейная задача наименьших квадратов}

Пусть дана переопределённая система
\[
    Ax\approx b,
    \qquad A\in\mathbb R^{m\times n},
    \qquad m\geq n.
\]
В общем случае точного решения нет. Метод наименьших квадратов выбирает
\[
    x_* = \arg\min_x \|Ax-b\|_2^2.
\]
Вектор
\[
    r=b-Ax_*
\]
называется невязкой.

Дифференцируя функционал
\[
    J(x)=(Ax-b)^T(Ax-b),
\]
получаем нормальные уравнения
\[
    A^TAx=A^Tb.
\]
Если столбцы $A$ линейно независимы, матрица $A^TA$ положительно определена и
решение единственно:
\[
    x_*=(A^TA)^{-1}A^Tb.
\]
На практике обратную матрицу явно не вычисляют.

\subsection{Геометрический смысл}

Вектор $Ax_*$ является ортогональной проекцией $b$ на столбцовое пространство
$\mathcal R(A)$. Поэтому
\[
    A^Tr=0.
\]
То есть невязка ортогональна каждому столбцу матрицы $A$. Если $A$ имеет
неполный столбцовый ранг, решений может быть несколько; среди них часто выбирают
решение минимальной нормы, задаваемое псевдообратной матрицей:
\[
    x_*=A^+b.
\]

\subsection{Численно устойчивое решение}

Нормальные уравнения ухудшают обусловленность:
\[
    \kappa_2(A^TA)=\kappa_2(A)^2.
\]
Поэтому предпочтительны:
\begin{enumerate}
    \item $QR$-разложение $A=QR$:
    \[
        \min_x\|Ax-b\|_2
        \quad\Longrightarrow\quad
        Rx=Q^Tb;
    \]
    \item сингулярное разложение
    \[
        A=U\Sigma V^T,
        \qquad
        A^+=V\Sigma^+U^T.
    \]
\end{enumerate}
SVD наиболее надёжно при вырожденности и позволяет определить численный ранг,
но дороже $QR$.

\subsection{Аппроксимация табличных данных}

Пусть заданы точки $(t_i,y_i)$ и модель
\[
    \varphi(t)=\sum_{j=0}^{n-1}c_j\phi_j(t).
\]
Минимизируется
\[
    \sum_{i=1}^{m}
    \left(y_i-\sum_{j=0}^{n-1}c_j\phi_j(t_i)\right)^2.
\]
Матрица задачи имеет элементы
\[
    A_{ij}=\phi_j(t_i).
\]
Для линейной регрессии $\varphi(t)=c_0+c_1t$ нормальные уравнения принимают вид
\[
\begin{pmatrix}
m&\sum t_i\\
\sum t_i&\sum t_i^2
\end{pmatrix}
\begin{pmatrix}c_0\\c_1\end{pmatrix}
=
\begin{pmatrix}\sum y_i\\\sum t_i y_i\end{pmatrix}.
\]

\subsection{Непрерывный метод наименьших квадратов}

В пространстве со скалярным произведением
\[
    (f,g)=\int_a^b w(x)f(x)g(x)\,dx,
    \qquad w(x)>0,
\]
ищется
\[
    \varphi_n=\sum_{j=0}^{n}c_j\phi_j
\]
из условия
\[
    \|f-\varphi_n\|_2^2\to\min.
\]
Нормальные уравнения:
\[
    \sum_{j=0}^{n}(\phi_j,\phi_i)c_j=(f,\phi_i),
    \qquad i=0,\ldots,n.
\]
Если базис ортогонален, система распадается:
\[
    c_i=\frac{(f,\phi_i)}{(\phi_i,\phi_i)}.
\]
При ортонормированном базисе $c_i=(f,\phi_i)$.

\subsection{Взвешенный и регуляризованный МНК}

Если наблюдения имеют различную точность, используют веса:
\[
    \min_x (Ax-b)^TW(Ax-b),
    \qquad W=W^T>0.
\]
Нормальные уравнения:
\[
    A^TWAx=A^TWb.
\]
При независимых ошибках с дисперсиями $\sigma_i^2$ обычно выбирают
$W_{ii}=1/\sigma_i^2$.

Для неустойчивых задач применяется регуляризация Тихонова:
\[
    \min_x\left(\|Ax-b\|_2^2+\lambda\|Lx\|_2^2\right),
\]
что приводит к системе
\[
    (A^TA+\lambda L^TL)x=A^Tb.
\]

\subsection{Нелинейный метод наименьших квадратов}

Для модели $r_i(x)$ минимизируется
\[
    J(x)=\frac12\sum_{i=1}^{m}r_i(x)^2.
\]
Метод Гаусса--Ньютона линеаризует невязку:
\[
    r(x+\Delta x)\approx r(x)+J_r(x)\Delta x
\]
и решает
\[
    J_r^TJ_r\Delta x=-J_r^Tr.
\]
Метод Левенберга--Марквардта добавляет стабилизирующий член:
\[
    (J_r^TJ_r+\lambda I)\Delta x=-J_r^Tr.
\]

\subsection{Статистическая интерпретация}

Если
\[
    b=Ax+\varepsilon,
    \qquad
    \varepsilon\sim\mathcal N(0,\sigma^2I),
\]
то при полном столбцовом ранге $A$ оценка МНК совпадает с оценкой максимального
правдоподобия. При стандартных условиях линейной модели она является несмещённой,
а её ковариационная матрица равна
\[
    \operatorname{Cov}(x_*)=\sigma^2(A^TA)^{-1}.
\]

\subsection{Что важно сказать на экзамене}

Нужно вывести нормальные уравнения, объяснить ортогональность невязки и отметить,
что вычислять решение через $(A^TA)^{-1}$ нежелательно. Численно надёжные методы
--- $QR$ и SVD. Важные расширения --- веса, регуляризация и нелинейный МНК.

\newpage